%=======================================================================
%  CHAPTER 4 — KNOWLEDGE REPRESENTATION, INFERENCE AND REASONING  (4 hrs)
%=======================================================================
\section{Knowledge Representation, Inference and Reasoning}

{\color{sub}\footnotesize 4 hrs \gap$\cdot$\gap asked in \mk{all 41 papers}
\gap$\cdot$\gap 4 syllabus sub-topics \gap$\cdot$\gap 638 marks, \mk{18.8\,\% of every
mark ever set}, against a 7.8\,\% share of the teaching hours \gap$\cdot$\gap
33 resolution proofs and 15 Bayes numericals worked in \emph{Ch4 -- Knowledge
Representation, Inference and Reasoning -- Solved Problems}.}

\lead{Read this first:}
\begin{itemize}
  \item \mk{This is the highest-paying chapter in the subject, by a factor
        of 2.4.} Four teaching hours, nearly a fifth of every mark. Whatever time the
        timetable suggests for it, spend more.
  \item \mk{One question dominates everything}: \textbf{prove a goal from
        given premises by resolution refutation} \tS{33}. It is in \textbf{33 of the 41
        papers}, always for 6 to 10 marks. The premises change; the procedure never does.
        Learn the procedure once and the question is answered for the rest of your life.
  \item Its two feeder topics are asked in their own right: \textbf{conversion to CNF}
        \tS{9} and the \textbf{rules of inference}. You cannot resolve without CNF, so
        this is one topic wearing two hats.
  \item The other two reliable topics: \textbf{forward Vs backward chaining} \tS{11},
        always a comparison with an example, and \textbf{belief / causal networks}
        \tF{7}.
  \item \mk{Fifteen papers set a Bayes numerical}, usually bolted onto a
        2--3 mark theory opener. The arithmetic is one formula; all 15 are worked in the
        companion.
  \item Pairs the examiner reliably asks together: a resolution proof preceded by
        \textbf{the CNF steps}, \textbf{skolemization} or \textbf{the rules of
        inference}, and a Bayes numerical preceded by \textbf{prior/posterior}. The short
        part is always the definition; the long part is always the work.
  \item \mk{The 5-mark saver}: if a resolution proof stalls in the exam,
        write the FOPL, write the clause form, write the negated goal, and show the
        resolution steps you do have. The marks are spread across those stages, not
        parked on the empty clause.
\end{itemize}

\hr

%-----------------------------------------------------------------------
\T{4.1 Knowledge, representation and reasoning}

\Q{What is knowledge, representation and reasoning? What kinds of knowledge are there?}{\m{2}\m{2+5}}
{\color{sub}\footnotesize \yr{\textbf{72 Ash} \m{2+5}, \textbf{77 Ch} \m{2}, 72 Ma \m{7}}
\gap$\cdot$\gap never asked alone: it is the opening 2 marks of a chaining or
rule-based-reasoning question, so \mk{answer it in four lines and move on}}

\sbs{%
\lead{The three words}
\begin{itemize}
  \item \textbf{Knowledge} $=$ facts, concepts, rules and relationships about a domain
        that an agent can use. It sits above data and information:
  \begin{itemize}
    \item \textbf{data} $\rightarrow$ raw symbols; \textbf{information} $\rightarrow$ data
          with context; \textbf{knowledge} $\rightarrow$ information plus how to use it;
          \textbf{intelligence} $\rightarrow$ knowledge applied to a new situation.
  \end{itemize}
  \item \textbf{Representation} $=$ writing that knowledge in a formal notation a machine
        can store and manipulate. Not any notation: one that
        \mk{supports inference}.
  \item \textbf{Reasoning} $=$ deriving new sentences from the stored ones, so the agent
        knows things nobody explicitly told it.
  \item The four questions KR exists to answer: how do we represent facts about the
        world, how do we reason about them, which representation suits the real world,
        and how do we put it in a form a machine can compute with.
\end{itemize}}{%
\lead{\mk{The five kinds of knowledge}}
\vspace{-1pt}
\begin{tabularx}{\linewidth}{@{}L{2.2cm}XL{2.6cm}@{}}
\toprule
\hrow \thead{Kind} & \thead{What it is} & \thead{Example} \\
\midrule
\textbf{Declarative} & Knowing \emph{that}. Facts, concepts, objects, stated without
saying how to use them & ``Kathmandu is the capital of Nepal'' \\
\textbf{Procedural} & Knowing \emph{how}. Rules, strategies, procedures, agendas &
the steps to sort a list \\
\textbf{Heuristic} & Rules of thumb from experience: good, not guaranteed &
in chess, keep the opponent's king exposed \\
\textbf{Structural} & Relationships between concepts and objects &
``a cow \emph{is-a} mammal'' \\
\textbf{Meta} & Knowledge \emph{about} knowledge: what the agent knows, and how reliable
it is & ``my map of Pokhara is out of date'' \\
\bottomrule
\end{tabularx}
\vspace{2pt}
\begin{itemize}
  \item \mk{Name all five with one example each.} Papers award a mark
        per kind, and \emph{meta-knowledge} and \emph{structural} are the two students
        forget.
\end{itemize}}

\penalty0\vspace{2pt}

\sbs{%
\lead{What a good representation must give you}
\begin{itemize}
  \item \textbf{Representational adequacy} $-$ can it express everything in the domain?
  \item \textbf{Inferential adequacy} $-$ can new knowledge be derived from it?
  \item \textbf{Inferential efficiency} $-$ can the derivation be steered, so the search
        does not explode?
  \item \textbf{Acquisitional efficiency} $-$ can new knowledge be added easily?
\end{itemize}}{%
\lead{Why logic is the representation of choice}
\begin{itemize}
  \item It is concise, unambiguous, context-insensitive and expressive, and it comes with
        a \textbf{proof theory}, i.e.\ \mk{inference is defined for you}.
        That last property is what the rest of this chapter is built on.
  \item Logic is fixed by three things, and papers ask for them by name:
        \textbf{syntax} (which strings are sentences), \textbf{semantics} (what a
        sentence means about the world) and \textbf{proof theory} (which new sentences
        may be derived from old ones).
\end{itemize}}

\penalty0\vspace{2pt}

\creamq{What is a knowledge-based agent? How does it work?}{\m{2+3+2}}
{\color{sub}\footnotesize \tP{1} \yr{\texttt{82 Ba} \m{2+3+2}} \gap$\cdot$\gap
\mk{no lecture deck in the set covers this by name} \added \gap$\cdot$\gap
answer it as the \emph{three operations} plus the two levels}

\sbs{%
\lead{Definition and structure}
\begin{itemize}
  \item An agent whose behaviour comes from a \textbf{knowledge base} (KB) of sentences
        about the world, plus an \textbf{inference engine} that derives new sentences
        from it. It \mk{acts on what it can prove}, not on hard-wired
        reflexes.
  \item Two levels, and papers like the distinction:
  \begin{itemize}
    \item \textbf{Knowledge level} $-$ \emph{what} the agent knows and what goal it has,
          independent of any notation.
    \item \textbf{Implementation level} $-$ the sentences, the data structures and the
          algorithm that actually run.
  \end{itemize}
  \item Three operations, in this order, once per time step:
  \begin{itemize}
    \item \textbf{TELL} the KB what was perceived;
    \item \textbf{ASK} the KB which action to take;
    \item \textbf{TELL} the KB that the action was taken.
  \end{itemize}
\end{itemize}}{%
\lead{Why it matters, with the standard example}
\begin{itemize}
  \item The \textbf{Wumpus world}: an agent in a cave perceives \emph{breeze},
        \emph{stench}, \emph{glitter}. No single percept names a pit, but
        \mk{breeze in (1,2) and (2,1) with none in (1,1)} \emph{proves}
        where the pit is. That inference, not the percept, is the whole point.
  \item \checkmark It can be \textbf{told new facts} without reprogramming: add a sentence
        and the behaviour changes.
  \item \checkmark It handles \textbf{partially observable} worlds, because it reasons
        about what it has not seen.
  \item \xmark It is only as good as its KB, and inference costs time: a large KB with an
        undirected engine is exactly the search-explosion problem of Ch\,3.
  \item \textbf{Every expert system} (Ch\,7) is a knowledge-based agent with a domain
        KB and a chaining inference engine. Say that sentence and the examiner sees you
        have connected the chapters.
\end{itemize}}

%-----------------------------------------------------------------------
\T{4.2 Propositional logic: WFF, tautology and the rules of inference}

\Q{Define a WFF with its properties. Prove that $(A \wedge (A \rightarrow B)) \rightarrow B$ is a tautology}{\m{4+4}}
{\color{sub}\footnotesize \tP{1} \yr{\textbf{\textit{71 Bh}} \m{4+4}} \gap$\cdot$\gap
the proof is a \mk{4-row truth table}, worked in the companion
\gap$\cdot$\gap ``short note: predicate logic / sentence validity'' \yr{\textbf{75 Bh}
\m{4}, \textit{72 Ma} \m{4}} wants the same vocabulary}

\sbs{%
\lead{Proposition and well-formed formula}
\begin{itemize}
  \item A \textbf{proposition} is a declarative statement that is \textbf{either true or
        false, never both}. ``5 is a prime number'' is one; ``close the door'' is not.
  \item An \textbf{atomic} sentence is a single propositional symbol ($P$, $Q$).
        A \textbf{compound} sentence joins atomic ones with the connectives
        $\neg, \wedge, \vee, \rightarrow, \leftrightarrow$.
  \item A \textbf{WFF} is a string built only by these recursive rules:
  \begin{itemize}
    \item any propositional symbol is a WFF;
    \item if $\alpha$ is a WFF then $\neg\alpha$ and $(\alpha)$ are WFFs;
    \item if $\alpha, \beta$ are WFFs then so are
          $(\alpha \wedge \beta)$, $(\alpha \vee \beta)$,
          $(\alpha \rightarrow \beta)$, $(\alpha \leftrightarrow \beta)$;
    \item \mk{nothing else is a WFF.} That closing clause is worth a
          mark and is the one students leave out.
  \end{itemize}
  \item Precedence, when brackets are missing:
        $\neg$, then $\wedge$, then $\vee$, then $\rightarrow$, then $\leftrightarrow$.
\end{itemize}}{%
\lead{\mk{The four properties to name}}
\begin{itemize}
  \item \textbf{Valid} (a \textbf{tautology}): true under \emph{every} interpretation.
        \emph{Eg} $P \vee \neg P$.
  \item \textbf{Satisfiable}: true under \emph{at least one}. \emph{Eg} $P \wedge Q$.
  \item \textbf{Unsatisfiable} (a \textbf{contradiction}): true under \emph{none}.
        \emph{Eg} $P \wedge \neg P$.
  \item \textbf{Equivalent}: $\alpha \equiv \beta$ when they have the same truth value in
        every row. Equivalently, $\alpha \leftrightarrow \beta$ is valid.
  \item \mk{The link resolution depends on}: $\alpha$ is valid
        \emph{iff} $\neg\alpha$ is unsatisfiable. That single line is why proving a goal
        is done by \textbf{refuting its negation} in \S4.5.
\end{itemize}
\lead{Equivalences you must be able to quote}
\begin{itemize}
  \item $\alpha \rightarrow \beta \;\equiv\; \neg\alpha \vee \beta$
        \gap\mk{(the workhorse)}
  \item $\alpha \leftrightarrow \beta \;\equiv\;
        (\alpha \rightarrow \beta) \wedge (\beta \rightarrow \alpha)$
  \item De Morgan: $\neg(\alpha \wedge \beta) \equiv \neg\alpha \vee \neg\beta$; \;
        $\neg(\alpha \vee \beta) \equiv \neg\alpha \wedge \neg\beta$
  \item Distribution:
        $\alpha \vee (\beta \wedge \gamma) \equiv
         (\alpha \vee \beta) \wedge (\alpha \vee \gamma)$
  \item \textbf{Contrapositive} $\neg\beta \rightarrow \neg\alpha$ is equivalent to
        $\alpha \rightarrow \beta$; the \textbf{converse}
        $\beta \rightarrow \alpha$ and the \textbf{inverse}
        $\neg\alpha \rightarrow \neg\beta$ are \textbf{not}.
\end{itemize}}

\penalty0\vspace{2pt}

\creamq{List the rules of inference. When can we use them?}{\m{2}\m{3}\m{2+3}}
{\color{sub}\footnotesize \yr{\textbf{78 Ch} \m{3}, 80 Ash \m{3}, \textbf{81 Ch} \m{2+3},
76 Ba \m{2}, 70 Ma \m{3}} \gap$\cdot$\gap 5 papers, none over 3 marks on its own, so no
tier chip \gap$\cdot$\gap \mk{it is always the opener to a resolution proof},
which is where the other 7 or 8 marks are}

\sbsr{0.56}{%
\lead{\mk{The eight propositional rules}}
\vspace{-1pt}
\begin{tabularx}{\linewidth}{@{}L{3.0cm}L{3.4cm}X@{}}
\toprule
\hrow \thead{Rule} & \thead{From \; $\therefore$ \; conclude} & \thead{Tautological form} \\
\midrule
\textbf{Modus ponens} & $P \rightarrow Q$, $P$ \; $\therefore$ \; $Q$ &
$[(P \rightarrow Q) \wedge P] \rightarrow Q$ \\
\textbf{Modus tollens} & $P \rightarrow Q$, $\neg Q$ \; $\therefore$ \; $\neg P$ &
$[(P \rightarrow Q) \wedge \neg Q] \rightarrow \neg P$ \\
\textbf{Hypothetical syllogism} & $P \rightarrow Q$, $Q \rightarrow R$ \; $\therefore$ \;
$P \rightarrow R$ & chains two implications \\
\textbf{Disjunctive syllogism} & $P \vee Q$, $\neg P$ \; $\therefore$ \; $Q$ &
eliminates a disjunct \\
\textbf{Addition} & $P$ \; $\therefore$ \; $P \vee Q$ & weakens \\
\textbf{Simplification} & $P \wedge Q$ \; $\therefore$ \; $P$ & splits \\
\textbf{Conjunction} & $P$, $Q$ \; $\therefore$ \; $P \wedge Q$ & joins \\
\textbf{Resolution} & $P \vee Q$, $\neg P \vee R$ \; $\therefore$ \; $Q \vee R$ &
\mk{the only one \S4.5 needs} \\
\bottomrule
\end{tabularx}}{%
\lead{The four quantifier rules}
\begin{itemize}
  \item \textbf{Universal instantiation}: from $\forall x\,P(x)$ infer $P(c)$ for
        \emph{any} constant $c$.
  \item \textbf{Existential instantiation}: from $\exists x\,P(x)$ infer $P(c)$ for a
        \mk{brand-new} constant $c$. Reusing an existing name is the
        classic mistake; this is skolemization in one line.
  \item \textbf{Universal generalisation}: from $P(c)$ for an arbitrary $c$ infer
        $\forall x\,P(x)$.
  \item \textbf{Existential generalisation}: from $P(c)$ infer $\exists x\,P(x)$.
\end{itemize}
\lead{\mk{When can we use them} (the second half of the question)}
\begin{itemize}
  \item Only when the premises \textbf{match the pattern exactly}, after unification.
        A rule is a \emph{syntactic} licence: it never looks at meaning.
  \item They are \textbf{sound}: applied to true premises they can only produce true
        conclusions, so a derivation is a proof.
  \item They are used when the KB is in the right form: modus ponens needs implications,
        \textbf{resolution needs CNF} (\S4.4), and forward/backward chaining
        (\S4.6) need \emph{definite} clauses.
\end{itemize}}

%-----------------------------------------------------------------------
\T{4.3 First-order predicate logic}

\Q{Why do we need FOPL? Explain its syntax and quantifiers}{\m{2}\m{2+3}\m{4}}
{\color{sub}\footnotesize \yr{70 Ma \m{2+3}, \textbf{75 Bh} \m{4}, \textit{72 Ma} \m{4}}
\gap$\cdot$\gap \mk{translating English into FOPL is the first step of every
resolution proof}, so this section is really 33 papers wide \gap$\cdot$\gap
a translation set is worked in the companion \yr{(\textbf{\textit{76 Bh}} \m{2})}}

\sbs{%
\lead{\mk{Why propositional logic is not enough}}
\begin{itemize}
  \item In PL a whole statement is \textbf{one indivisible symbol}. ``All humans are
        intelligent'' and ``Some humans are intelligent'' become two unrelated letters
        $P$ and $Q$, with no way to see that one implies the other.
  \item It cannot express \textbf{objects, their properties or relations between them},
        and it has \mk{no quantifiers}, so a fact about a million
        individuals needs a million sentences.
  \item FOPL fixes exactly that: it breaks a sentence into \textbf{subject} and
        \textbf{predicate}, and quantifies over objects.
  \item The three-line answer: PL has \mk{limited expressive power};
        FOPL adds objects, predicates and quantifiers; so one FOPL sentence replaces an
        unbounded set of PL ones.
\end{itemize}
\lead{Syntax: the seven elements}
\begin{itemize}
  \item \textbf{Constants} $-$ \emph{John}, \emph{Nepal}, 2. Named objects.
  \item \textbf{Variables} $-$ $x, y, z$.
  \item \textbf{Predicates} $-$ \emph{Brother}, \emph{Loves}, $>$. Relations, true or
        false of their arguments.
  \item \textbf{Functions} $-$ \emph{FatherOf}, \emph{sqrt}. Return an \emph{object},
        not a truth value. \mk{Predicate Vs function is a favourite
        one-mark trap.}
  \item \textbf{Connectives} $\neg, \wedge, \vee, \rightarrow, \leftrightarrow$,
        \textbf{equality} $=$, and \textbf{quantifiers} $\forall, \exists$.
\end{itemize}}{%
\lead{The two quantifiers, and the connective each takes}
\begin{itemize}
  \item $\forall$ (\textbf{universal}, ``for all'') goes with $\rightarrow$:\\
        \mbox{}\hfill $\forall x\; man(x) \rightarrow drinks(x, coffee)$\hfill\mbox{}
  \item $\exists$ (\textbf{existential}, ``there exists'') goes with $\wedge$:\\
        \mbox{}\hfill $\exists x\; student(x) \wedge fails(x, maths)$\hfill\mbox{}
  \item \mk{Swap the connectives and the sentence is nonsense.}
        $\forall x\, (man(x) \wedge drinks(x,coffee))$ says \emph{everything in the
        universe is a coffee-drinking man}; $\exists x\, (student(x) \rightarrow
        fails(x))$ is satisfied by any non-student. This single pairing earns or loses
        the translation marks.
  \item \textbf{Duality}, needed in every CNF conversion:
        $\neg \forall x\, P \equiv \exists x\, \neg P$ \; and \;
        $\neg \exists x\, P \equiv \forall x\, \neg P$.
\end{itemize}
\lead{Translation patterns worth memorising}
\begin{itemize}
  \item ``All A are B'' $\rightarrow$ $\forall x\, A(x) \rightarrow B(x)$
  \item ``Some A are B'' $\rightarrow$ $\exists x\, A(x) \wedge B(x)$
  \item ``No A is B'' $\rightarrow$ $\forall x\, A(x) \rightarrow \neg B(x)$, or
        $\neg \exists x\, (A(x) \wedge B(x))$
  \item ``A has at least one B'' $\rightarrow$ $\exists x\, B(x) \wedge has(A, x)$
  \item ``Everyone who \ldots\ all animals'' $\rightarrow$ a $\forall$ \emph{inside} the
        antecedent of another $\rightarrow$; that nesting is what makes the standard CNF
        example in \S4.4 hard.
\end{itemize}}

\penalty0\vspace{2pt}

\sbs{%
\lead{Skolemization \yr{(\texttt{81 Ba} \m{2}, \textbf{69 Bh}, short note)}}
\begin{itemize}
  \item \textbf{Definition}: removing an existential quantifier by replacing its variable
        with a \emph{new} symbol, so the sentence can be written as a clause.
  \item \textbf{Two cases, and the second is the one papers test}:
  \begin{itemize}
    \item $\exists$ \textbf{not inside any} $\forall$ $\rightarrow$ replace with a new
          \textbf{Skolem constant}. \; $\exists x\, rich(x) \Rightarrow rich(A)$.
    \item $\exists$ \textbf{inside} $\forall$ $\rightarrow$ replace with a
          \textbf{Skolem function} of the enclosing universal variables. \;
          $\forall x\, \exists y\, loves(x,y) \Rightarrow loves(x, F(x))$.
  \end{itemize}
  \item \mk{Why a function and not a constant}: the person loved
        depends on who is loving. A constant would assert that
        \emph{everyone loves the same individual}, which is a different and stronger
        claim, and the proof built on it is wrong.
  \item The new symbol must be one \textbf{not already used} anywhere in the KB.
\end{itemize}}{%
\lead{Horn clause and definite clause \yr{(\textbf{78 Ch} \m{3}, 70 Ma)}}
\begin{itemize}
  \item A \textbf{clause} is a disjunction of literals: $\neg P \vee Q \vee \neg R$.
  \item \textbf{Definite clause}: a clause with \mk{exactly one
        positive} literal. \; $\neg P \vee \neg Q \vee K$.
  \item \textbf{Horn clause}: a clause with \mk{at most one positive}
        literal. So every definite clause is a Horn clause, and the extra case is the
        \textbf{goal clause} with no positive literal at all.
  \item \textbf{Why the restriction exists}, which is what the 3 marks are for: a definite
        clause is exactly an implication with a single conclusion,
        $(P \wedge Q) \rightarrow K$, so
  \begin{itemize}
    \item inference collapses to \textbf{modus ponens}, and
    \item \textbf{forward and backward chaining} (\S4.6) become possible, in time
          \emph{linear} in the size of the KB rather than exponential.
  \end{itemize}
  \item \textbf{Prolog} is exactly this: a program is a set of definite clauses.
\end{itemize}}

\penalty0\vspace{2pt}

\creamq{Unification and the most general unifier}{\m{2}\m{3}}
{\color{sub}\footnotesize \yr{asked inside resolution questions, never alone}
\gap$\cdot$\gap \mk{but no resolution step is legal without it}, and
examiners take marks off a proof that unifies wrongly}

\sbs{%
\lead{What it does}
\begin{itemize}
  \item \textbf{Unification} finds a substitution $\theta$ that makes two atomic
        sentences identical: $\text{UNIFY}(\alpha, \beta) = \theta$ where
        $\alpha\theta = \beta\theta$.
  \item The substitution with the fewest commitments is the \textbf{most general unifier
        (MGU)}, and it is the one to use: a more specific one can lose proofs.
  \item Written $[term/variable]$, \emph{Eg} $\theta = [John/x]$.
  \item \textbf{It fails}, and the resolution step is then illegal, when
  \begin{itemize}
    \item the predicate symbols or the arities differ;
    \item two different constants must match, \emph{Eg} $King(John)$ Vs $King(Ram)$;
    \item a variable would have to be bound to a term containing itself
          (the \emph{occurs check}).
  \end{itemize}
\end{itemize}}{%
\lead{Worked, in the form the papers expect}
\begin{itemize}
  \item $King(x)$ and $King(John)$ \; $\rightarrow$ \; $\theta = [John/x]$.
  \item $P(x, y)$ and $P(a, f(z))$ \; $\rightarrow$ \; $\theta = [a/x,\; f(z)/y]$.
  \item $P(x, f(y))$ and $P(a, f(g(x)))$ \; $\rightarrow$ \; $\theta = [a/x,\; g(a)/y]$:
        substitute $a$ for $x$ first, then \mk{apply that binding inside
        the second term too}. Writing $g(x)/y$ and stopping is the common
        half-answer.
  \item $Q(a, g(x,a), f(y))$ and $Q(a, g(f(b),a), x)$ \; $\rightarrow$ \;
        $\theta = [f(b)/x,\; b/y]$.
  \item \textbf{Standardise apart first}: if both clauses use $x$, rename one of them.
        Two clauses that share a variable name will unify wrongly and the proof will be
        rejected.
\end{itemize}}

%-----------------------------------------------------------------------
\T{4.4 Conversion to conjunctive normal form}

\Q{Why is CNF required? Explain all the steps involved in converting to CNF, with an example}{\m{2+6}\m{3}\m{5}\m{8}\m{10}}
{\color{sub}\footnotesize \tS{9} \yr{69 Po \m{10}, \textbf{70 Bh} \m{8}, 75 Ba \m{2+6},
\textbf{79 Ch} \m{2+6}, \textbf{75 Bh} \m{2+6}, \textbf{\textit{76 Bh}} \m{5},
\textbf{80 Ch} \m{3}, \textbf{81 Ch} \m{2+3}, 76 Ba \m{4}} \gap$\cdot$\gap
\mk{four more papers} give a formula and ask for its CNF \tF{4}
\yr{\textbf{\textit{73 Bh}} \m{2+2+2+3}, \textit{74 Ma} \m{6},
\textbf{\textit{72 Ash}} \m{8}, 79 Jth \m{2+2+4}}, all worked in the companion}

\sbs{%
\lead{\mk{Why CNF is required} (the 2 marks, and the only part
students skip)}
\begin{itemize}
  \item \textbf{CNF} $=$ a conjunction of clauses, each clause a disjunction of literals:
        \mbox{}\hfill $(A \vee \neg B) \wedge (C \vee D \vee \neg E)$ \hfill\mbox{}
  \item \textbf{Resolution is a single inference rule and it only operates on clauses.}
        Everything else $-$ implications, biconditionals, nested negations, quantifiers
        $-$ has to go before the machinery can start.
  \item Because the whole formula is one big \textbf{AND}, each clause can be
        \mk{asserted separately} and stored as its own line in the KB.
  \item \textbf{Every} sentence of propositional or first-order logic has a CNF
        equivalent, so nothing is lost by insisting on it. That completeness is what
        makes resolution a \emph{general} proof procedure.
  \item Contrast \textbf{DNF}, a disjunction of conjunctions. Papers sometimes ask for
        both: CNF is joined by AND, DNF by OR.
\end{itemize}}{%
\lead{\mk{The eight steps, in this order}}
\begin{enumerate}
  \item \textbf{Eliminate} $\leftrightarrow$: \;
        $\alpha \leftrightarrow \beta \Rightarrow
        (\alpha \rightarrow \beta) \wedge (\beta \rightarrow \alpha)$.
  \item \textbf{Eliminate} $\rightarrow$: \;
        $\alpha \rightarrow \beta \Rightarrow \neg\alpha \vee \beta$.
  \item \textbf{Move $\neg$ inwards}, until it sits on atoms only: De Morgan,
        $\neg\neg\alpha \equiv \alpha$, and the quantifier duals
        $\neg\forall x\,P \equiv \exists x\,\neg P$,
        $\neg\exists x\,P \equiv \forall x\,\neg P$.
  \item \textbf{Standardise variables apart}: rename so that no two quantifiers bind the
        same variable name.
  \item \textbf{Skolemize}: drop each $\exists$, replacing its variable by a new constant,
        or by a \textbf{function} of the enclosing universal variables (\S4.3).
  \item \textbf{Drop the universal quantifiers}: every remaining variable is understood
        to be universally quantified.
  \item \textbf{Distribute $\vee$ over $\wedge$}: \;
        $\alpha \vee (\beta \wedge \gamma) \Rightarrow
        (\alpha \vee \beta) \wedge (\alpha \vee \gamma)$.
  \item \textbf{Split} the conjunction into separate clauses, and rename variables again
        so no two clauses share one.
\end{enumerate}
\vspace{1pt}
{\color{sub}\footnotesize \mk{Steps 4, 7 and 8 are the ones the local
textbook leaves out}, and a proof that skips 4 or 8 unifies two unrelated
variables and derives nonsense. Write all eight; the marks are per step.}}

\penalty0\vspace{2pt}

\sbs{%
\lead{The worked example the papers use
\yr{(\textbf{79 Ch}, \textbf{75 Bh}, \textbf{\textit{73 Bh}})}}
\begin{itemize}
  \item \emph{``Everyone who loves all animals is loved by someone.''}
  \item \textbf{FOPL}: \;
        $\forall x\,[\forall y\; Animal(y) \rightarrow Loves(x,y)]
         \rightarrow [\exists z\; Loves(z,x)]$
  \item \textbf{2. Eliminate} $\rightarrow$, twice: \;
        $\forall x\, \neg[\forall y\; \neg Animal(y) \vee Loves(x,y)]
         \vee [\exists z\; Loves(z,x)]$
  \item \textbf{3. Move $\neg$ in}, using the dual and De Morgan: \;
        $\forall x\, [\exists y\; Animal(y) \wedge \neg Loves(x,y)]
         \vee [\exists z\; Loves(z,x)]$
  \item \textbf{5. Skolemize.} Both $\exists$ are inside $\forall x$, so both become
        \textbf{functions of $x$}: $y \Rightarrow F(x)$, $z \Rightarrow G(x)$: \;
        $\forall x\, [Animal(F(x)) \wedge \neg Loves(x, F(x))] \vee Loves(G(x), x)$
  \item \textbf{6. Drop $\forall$}, \textbf{7. distribute}, \textbf{8. split}:
  \begin{itemize}
    \item $Animal(F(x)) \vee Loves(G(x), x)$
    \item $\neg Loves(x, F(x)) \vee Loves(G(x), x)$
  \end{itemize}
  \item \mk{Read the Skolem functions back}: $F(x)$ is ``the animal
        $x$ does not love, if there is one'' and $G(x)$ is ``someone who loves $x$''.
        Saying that in one line shows the step was understood, not copied.
\end{itemize}}{%
\lead{\mk{Where marks are actually lost}}
\begin{itemize}
  \item \textbf{Skolemizing to a constant inside a $\forall$.} $\forall x \exists y$
        becomes $F(x)$, never $A$. This is the single commonest error, and it silently
        changes the meaning.
  \item \textbf{Distributing where you must not.}
        $\neg(A \wedge \neg B) \vee C$ is
        $\neg A \vee B \vee C$ $-$ \emph{one} clause. It is \textbf{not}
        $(\neg A \vee C) \wedge (B \vee C)$.
        \mk{The local textbook makes exactly this error} in its
        ``likes all food'' problem, then quietly uses the correct single clause in the
        proof underneath. The companion works that family properly.
  \item \textbf{Forgetting to standardise apart.} Two clauses both using $x$ will unify on
        a variable they do not share in meaning.
  \item \textbf{Dropping $\forall$ before skolemizing.} The order in the list is not
        decoration: once the $\forall$ is gone there is nothing to make the Skolem
        function depend on.
  \item \textbf{Simplification is allowed and is worth doing}: a clause containing both
        $P$ and $\neg P$ is a tautology and can be dropped, and duplicate literals in one
        clause collapse.
\end{itemize}}

%-----------------------------------------------------------------------
\T{4.5 Resolution refutation}

\Q{Prove the given goal from the given premises using resolution refutation}{\m{6}\m{7}\m{8}\m{9}\m{10}}
{\color{sub}\footnotesize \tS{33} \yr{\textbf{69 Bh} \m{10}, 75 Ba \m{8},
\textbf{77 Ch} \m{8}, 82 Ka \m{8}, 81 Ash \m{2+6}, \textbf{\texttt{81 Bh}} \m{9},
\textbf{\texttt{82 Bh}} \m{6}, \textbf{70 Bh} \m{8}, \textbf{71 Bh} \m{8},
\textbf{72 Ash} \m{7}, \textbf{74 Bh} \m{8}, \textbf{\textit{73 Bh}} \m{8},
\textbf{\textit{71 Bh}} \m{8}, \textbf{\texttt{80 Bh}} \m{8}, \texttt{80 Ba} \m{8},
\textit{72 Ma} \m{6+4}, 72 Ma \m{7}, \textbf{73 Bh} \m{8}, 78 Po \m{8},
\textbf{\textit{74 Bh}} \m{8}, 79 Jth \m{6}, \textbf{\texttt{79 Bh}} \m{8}, 78 Ba \m{8},
73 Ma \m{10}, \texttt{81 Ba} \m{2+6}, \textbf{81 Ch} \m{8}, \textbf{78 Ch} \m{3+3},
\textbf{80 Ch} \m{3+5}, 80 Ash \m{3+5}, 79 Ash \m{7}, 76 Ba \m{2+8},
\textbf{\textit{76 Bh}} \m{7}, \textbf{\textit{77 Ch}} \m{8}} \gap$\cdot$\gap
\mk{33 of 41 papers, 260 marks} \gap$\cdot$\gap
every one of them is worked in the companion}

\sbsr{0.52}{%
\lead{\mk{The procedure. This is the chapter.}}
\begin{itemize}
  \item \textbf{1.} Convert every premise into \textbf{FOPL}.
  \item \textbf{2.} Convert every FOPL sentence into \textbf{CNF} (\S4.4), and number the
        clauses.
  \item \textbf{3.} \mk{Negate the goal}, convert that to CNF too, and
        add it to the clause set.
  \item \textbf{4.} Repeatedly \textbf{resolve} two clauses containing complementary
        literals, unifying them with the MGU, and add the \textbf{resolvent}.
  \item \textbf{5.} Stop when the \textbf{empty clause} $\square$ is derived. The clause
        set is then unsatisfiable, so the negated goal is false, so
        \mk{the goal is true}.
\end{itemize}
\lead{Why refute instead of prove directly}
\begin{itemize}
  \item Resolution is \textbf{refutation complete}, not complete: it is guaranteed to find
        a contradiction in an unsatisfiable set, but it is not guaranteed to derive an
        arbitrary true sentence.
  \item So the proof is turned round using \; $KB \models \alpha$ \; iff \;
        $KB \wedge \neg\alpha$ \; is unsatisfiable, the equivalence from \S4.2.
\end{itemize}}{%
\lead{\mk{How to write one that scores}}
\begin{itemize}
  \item \textbf{Number every clause} and, beside each resolvent, write \emph{which two
        clauses} it came from and \emph{which substitution} was used. Examiners award
        those annotations.
  \item \textbf{Resolve towards the goal.} Start from the negated goal clause and work
        down: it prunes the search and it reads as a proof rather than a search log.
  \item \textbf{One complementary pair per step.} Cancelling two pairs in one line is
        unsound and is a standard trap.
  \item \textbf{Standardise variables apart} before each step (\S4.4, step 8).
  \item Finish with the sentence \emph{``the empty clause is derived, so the assumption
        contradicts the premises, hence the goal is proved''}. It is worth a mark and
        takes one line.
\end{itemize}}

\penalty0\vspace{2pt}

\sbsr{0.56}{%
\figT{c4_res_graph.png}
{\color{sub}\footnotesize The shape the answer must take: clauses down the left, each
resolvent to its right, the empty clause $\square$ at the foot. Premises are the Colonel
West set, asked in 7 papers.}}{%
\lead{Contradiction Vs extraction \yr{(the textbook's two methods)}}
\begin{itemize}
  \item \textbf{Answer by contradiction}: negate the goal, derive $\square$. Use this
        whenever the paper says \emph{prove that \ldots}, which is 21 of the 23 sets.
  \item \textbf{Answer by extraction}: do not negate; chain the clauses forward until the
        goal literal itself is produced. Use it when the paper asks \emph{who} or
        \emph{what} rather than \emph{prove}, \emph{Eg} \yr{72 Ma}'s
        \emph{``did Marcus hate Caesar?''}
  \item \mk{Same clauses, same resolution rule}; only the starting
        point differs.
\end{itemize}}

\penalty0\vspace{2pt}

\sbs{%
\lead{\mk{33 papers, but only 26 different premise sets}, and four
families cover 22 of them}
\begin{itemize}
  \item \textbf{Likes-all-food / peanuts} $-$ \textbf{7 papers}, five wordings of one
        puzzle (John/Bill/Sue, Ram/Krishna/Radha, Ram/Anil, Sita/mushroom,
        Bhogendra/chicken). \mk{Learn this one first.}
  \item \textbf{Charlie is a horse / a mammal} $-$ 6 papers, three wordings. Tests
        \emph{inverse relations}, and \mk{two of its six premises are never
        used}.
  \item \textbf{Selling weapons to a hostile nation} $-$ 7 papers under six names, one
        asking for forward chaining instead. Tests \emph{existential instantiation}.
  \item \textbf{Light sleeper, hounds and mice} $-$ 2 papers. The hardest of the four: the
        goal is itself an implication, so \mk{negating it gives three
        clauses}.
  \item Then 10 one-off sets, from Marcus and Caesar to Curiosity and the cat.
\end{itemize}}{%
\lead{\mk{The two things that go wrong in an exam}}
\begin{itemize}
  \item \textbf{Negating a goal that is an implication.} \emph{``If John is a light
        sleeper then John has no mice''} negates to \emph{``John IS a light sleeper AND
        John HAS a mouse''}, which after skolemization is
        \textbf{three} clauses: $LS(John)$, $Have(John, b)$, $Mouse(b)$. Students negate
        it to one clause and the proof cannot close.
  \item \textbf{``Has some'' is existential.} ``Nono has some missiles'' is
        $\exists x\, Missile(x) \wedge Owns(Nono, x)$, which skolemizes to a
        \textbf{new constant} $M_1$: two clauses, $Missile(M_1)$ and $Owns(Nono, M_1)$.
        Writing $Missile(x)$ with a free variable claims \emph{everything} is a missile.
  \item Both are worked step by step in the companion, which prints each paper's premises
        exactly as the paper words them.
\end{itemize}}

%-----------------------------------------------------------------------
\T{4.6 Rule-based reasoning: forward and backward chaining}

\Q{Differentiate between forward and backward chaining with a suitable example}{\m{2}\m{3}\m{4}\m{7}\m{8}}
{\color{sub}\footnotesize \tS{11} \yr{70 Ma \m{4+4}, \textbf{74 Bh} \m{4},
\textbf{\texttt{79 Bh}} \m{7}, \textbf{\texttt{81 Bh}} \m{7}, 82 Ka \m{2+4},
81 Ash \m{2}, \textbf{\textit{73 Bh}} \m{3}, \textbf{80 Ch} \m{4},
\textbf{72 Ash} \m{2+5}, 78 Po \m{8}, \textbf{\textit{77 Ch}} \m{8}} \gap$\cdot$\gap
\mk{always wants an example, not just the table} \gap$\cdot$\gap
in 81 Ash the example \emph{is} the Colonel West proof, done forward}

\sbs{%
\lead{Forward chaining: data-driven}
\begin{itemize}
  \item Start from the \textbf{known facts}. Fire every rule whose premises are all
        satisfied, add its conclusion to the fact base, repeat until the goal appears or
        nothing new can be derived.
  \item \textbf{Bottom-up}, and it uses \textbf{modus ponens} on definite clauses.
  \item \checkmark Derives \textbf{everything} that follows, so it suits monitoring,
        alarms and planning, where new data arrives and consequences must be found.
  \item \xmark Does a lot of \textbf{irrelevant work}: it proves facts nobody asked for.
\end{itemize}}{%
\lead{Backward chaining: goal-driven}
\begin{itemize}
  \item Start from the \textbf{goal}. Find a rule that concludes it, make that rule's
        premises the new sub-goals, and recurse until every sub-goal is a known fact.
  \item \textbf{Top-down}, \textbf{depth first}, and again modus ponens.
  \item \checkmark Touches \textbf{only what the goal needs}, so it suits diagnosis, query
        answering and theorem proving. \textbf{Prolog} works this way.
  \item \xmark Can loop on recursive rules, and it answers only the question asked.
\end{itemize}}

\penalty0\vspace{2pt}

\sbs{%
\figT{c4_fwd_chain.png}
{\color{sub}\footnotesize \textbf{Forward} chaining on the Colonel West premises: facts
on the bottom row, one new row per iteration, until $Criminal(Robert)$ appears.}}{%
\figT{c4_bwd_chain.png}
{\color{sub}\footnotesize \textbf{Backward} chaining, same premises: the goal at the top,
its four conjuncts below it, each expanded until it meets a fact. The binding that made
each unification succeed sits under its node.}}

\penalty0\vspace{2pt}

\begin{tabularx}{\textwidth}{@{}L{2.9cm}XX@{}}
\toprule
\hrow & \thead{Forward chaining} & \thead{Backward chaining} \\
\midrule
Starts from & Known facts / data & The goal to be proved \\
Direction & Bottom-up, data-driven & Top-down, goal-driven \\
Question answered & ``What can I conclude?'' & ``Can I prove this?'' \\
Search & Breadth-first in effect: all consequences of the current facts & Depth-first: one sub-goal chain at a time \\
Work done & Derives many facts, most irrelevant to any one goal & Derives only what the goal needs \\
Suits & Monitoring, alarms, planning, real-time control, design & Diagnosis, query answering, theorem proving, Prolog \\
Number of goals & Handles \textbf{many} goals at once & Handles \textbf{one} goal at a time \\
Efficiency & Slow when the goal is specific & Slow when the fact base is small and rules are deep \\
\bottomrule
\end{tabularx}
\par\penalty0
\vspace{3pt}

\begin{itemize}
  \item \mk{The one-line discriminator to write}: forward chaining
        asks \emph{what follows from what I know}, backward chaining asks
        \emph{what would I need to know for this to be true}. Same rules, same modus
        ponens, opposite ends.
\end{itemize}

\penalty0\vspace{2pt}

\creamq{What is rule-based reasoning?}{\m{2}\m{7}}
{\color{sub}\footnotesize \tF{4} \yr{72 Ma \m{7}, 70 Ma \m{4+4}, \textbf{71 Bh} \m{7},
\textbf{77 Ch} \m{2}} \gap$\cdot$\gap asked as the \mk{2-mark opener} to
a backward-chaining question in three of the four}

\sbs{%
\lead{Definition and parts}
\begin{itemize}
  \item Knowledge stored as \textbf{IF \emph{condition} THEN \emph{action}} rules, and
        reasoning done by an \textbf{inference engine} that decides which rule to fire.
  \item Three components:
  \begin{itemize}
    \item \textbf{Rule base} $-$ the IF-THEN rules (long-term knowledge);
    \item \textbf{Working memory} $-$ the facts known right now;
    \item \textbf{Inference engine} $-$ match, resolve conflicts, fire.
  \end{itemize}
  \item The engine's cycle is \textbf{match $\rightarrow$ conflict resolution
        $\rightarrow$ execute}, exactly the production system of Ch\,2.
        \mk{Say so}: the examiner set that question too.
\end{itemize}}{%
\lead{Why it is used, and the worked shape}
\begin{itemize}
  \item \checkmark Rules are \textbf{modular}: add or remove one without touching the
        rest. \checkmark They are readable by a domain expert who is not a programmer.
        \checkmark The chain of fired rules \textbf{is} the explanation, which is how an
        expert system justifies itself (Ch\,7).
  \item \xmark No learning, and \xmark large rule bases become slow and can conflict.
  \item \textbf{The example to give}, since every paper asks for one:
  \begin{itemize}
    \item R1: IF \emph{fever} AND \emph{cough} THEN \emph{flu}.
          R2: IF \emph{flu} THEN \emph{prescribe rest}.
    \item Facts: \emph{fever}, \emph{cough}. Forward: R1 fires, \emph{flu} added, R2
          fires, \emph{prescribe rest}. Backward from \emph{prescribe rest}: needs
          \emph{flu}, which needs \emph{fever} and \emph{cough}, both known.
          \mk{Two directions, one rule base.}
  \end{itemize}
\end{itemize}}

%-----------------------------------------------------------------------
\T{4.7 Statistical reasoning: uncertainty, Bayes and belief networks}

\Q{How is reasoning done under uncertainty? What are the approaches for measuring it?}{\m{2}\m{3}\m{3+6}}
{\color{sub}\footnotesize \yr{80 Ash \m{3}, 81 Ash \m{2}, 79 Ash \m{3+6}}
\gap$\cdot$\gap \mk{always the opener to a Bayes numerical}, so keep it to
half a page and spend the time on the calculation}

\sbs{%
\lead{Why logic is not enough}
\begin{itemize}
  \item FOPL is \textbf{monotonic and categorical}: a fact is true or false, and adding
        knowledge never retracts a conclusion. The real world is neither.
  \item Three sources of uncertainty, and papers want them named:
  \begin{itemize}
    \item \textbf{incomplete} knowledge (not all facts are available);
    \item \textbf{imprecise / unreliable} data (a sensor, a symptom);
    \item \textbf{rules that are only usually true} (\emph{a stiff neck suggests
          meningitis}).
  \end{itemize}
  \item Writing \emph{stiff neck $\rightarrow$ meningitis} in logic is simply false, and
        listing every exception is impossible. \mk{So attach a number to
        the belief instead.}
\end{itemize}}{%
\lead{\mk{The approaches for measuring uncertainty}}
\begin{itemize}
  \item \textbf{Probability theory} $-$ degree of belief in $[0,1]$, updated by
        \textbf{Bayes' rule}. The mainstream answer, and the only one the numericals use.
  \item \textbf{Certainty factors} $-$ an ad-hoc $[-1, +1]$ score attached to each rule,
        as used in MYCIN. Simple, but not probabilistically sound.
  \item \textbf{Dempster-Shafer theory} $-$ assigns belief to \emph{sets} of hypotheses,
        so it can represent \emph{ignorance} separately from \emph{disbelief}.
  \item \textbf{Fuzzy logic} $-$ handles \emph{vagueness} (``tall'', ``warm''), which is a
        different thing from uncertainty. Covered in Ch\,6.
  \item \textbf{Non-monotonic / default reasoning} $-$ assume the normal case, retract the
        conclusion when contradicted.
  \item \mk{The distinction worth one mark}: probability is about
        \emph{not knowing which case holds}; fuzzy logic is about
        \emph{the case not having a sharp boundary}.
\end{itemize}}

\penalty0\vspace{2pt}

\Q{What is prior and posterior probability? Explain Bayes' theorem and its role in reasoning}{\m{2}\m{3}\m{2+6}\m{3+5}}
{\color{sub}\footnotesize \tS{15} \yr{71 Ma \m{3+5}, \textbf{73 Bh} \m{7},
73 Ma \m{8}, \textbf{\textit{73 Bh}} \m{8}, \textbf{\textit{72 Ash}} \m{8},
\textbf{78 Ch} \m{3+5}, \textbf{80 Ch} \m{3+5}, \textbf{\texttt{80 Bh}} \m{2+6},
\texttt{80 Ba} \m{5+3}, \texttt{81 Ba} \m{2+6}, 80 Ash \m{3+5}, 81 Ash \m{2+6},
78 Ba \m{2+6}, \texttt{82 Ba} \m{2+6}, \textbf{\texttt{82 Bh}} \m{7}} \gap$\cdot$\gap
\mk{15 papers carry a Bayes calculation}, every one of them worked in the
companion \gap$\cdot$\gap the theory part is never more than 3 marks}

\sbsr{0.46}{%
\lead{The vocabulary}
\begin{itemize}
  \item \textbf{Prior} $P(A)$ $-$ belief in $A$ \mk{before} any
        evidence. \emph{Eg} 1 person in 50,000 has meningitis.
  \item \textbf{Posterior} $P(A|B)$ $-$ belief in $A$ \mk{after}
        evidence $B$ is observed. \emph{Eg} the chance of meningitis \emph{given} a stiff
        neck.
  \item \textbf{Likelihood} $P(B|A)$ $-$ how probable the evidence is if the hypothesis
        holds. \mk{This is the number a doctor actually knows}, and it is
        the reason the theorem is needed.
  \item \textbf{Marginal} $P(B)$ $-$ the probability of the evidence however it arose.
\end{itemize}
\lead{The theorem}
\[ P(A|B) \;=\; \frac{P(B|A)\,P(A)}{P(B)} \]
\begin{itemize}
  \item When $P(B)$ is not given, expand it by the
        \textbf{law of total probability}, which is what most papers require:
\end{itemize}
\[ P(A|B) = \frac{P(B|A)P(A)}{P(B|A)P(A) + P(B|\neg A)P(\neg A)} \]}{%
\lead{\mk{How it aids inference} (the marks in the theory half)}
\begin{itemize}
  \item It \mk{reverses the direction of a conditional}. Causal
        knowledge runs cause $\rightarrow$ effect ($P(\textit{symptom}\,|\,
        \textit{disease})$, stable and measurable); diagnosis needs effect
        $\rightarrow$ cause. Bayes converts one into the other.
  \item It gives a \textbf{principled update rule}: posterior $\propto$ likelihood
        $\times$ prior. Evidence does not replace belief, it revises it.
  \item It makes the \textbf{base rate} explicit. A 99\,\%-accurate test for a
        1-in-10,000 disease still leaves you far more likely to be healthy than ill,
        because the prior is tiny. \mk{This is asked directly}
        \yr{(\textbf{\textit{72 Ash}})} and it is the single most useful sentence in the
        section.
  \item It composes: a chain of evidence can be applied one piece at a time, which is
        exactly what a belief network does below.
\end{itemize}}

\penalty0\vspace{2pt}

\sbsr{0.62}{%
\figT{c4_bayes_anatomy.png}}{%
\lead{Getting the arithmetic right}
\begin{itemize}
  \item Name the events \textbf{before} substituting. Half the lost marks are a $P(A|B)$
        written where $P(B|A)$ was meant.
  \item If the question gives percentages of a population that split into groups
        (\emph{factories}, \emph{men and women}), the denominator is
        \mk{always the total-probability sum} over those groups.
  \item $P(\neg A) = 1 - P(A)$ is almost always needed and almost never given.
\end{itemize}}

\penalty0\vspace{2pt}

\Q{What is a causal / belief network? How is reasoning done in it?}{\m{2+5}\m{2+6}\m{6}\m{7}}
{\color{sub}\footnotesize \tF{7} \yr{71 Ma \m{2+5}, \textbf{79 Ch} \m{2+6},
79 Ash \m{3+6}, \textbf{77 Ch} \m{6}, \textbf{73 Bh} \m{7}, \texttt{80 Ba} \m{5+3},
\textbf{\texttt{82 Bh}} \m{7}} \gap$\cdot$\gap
\emph{causal net}, \emph{belief network}, \emph{Bayesian network} and \emph{Bayesian
belief network} are \mk{four names for one thing} in these papers
\gap$\cdot$\gap two papers print a network and ask for a joint probability
\yr{(78 Ba, \textbf{\texttt{82 Bh}})}}

\sbsr{0.54}{%
\lead{Definition and the two parts}
\begin{itemize}
  \item A \textbf{directed acyclic graph} in which
  \begin{itemize}
    \item each \textbf{node} is a random variable, and
    \item each \textbf{edge} $X \rightarrow Y$ says $X$ has a \textbf{direct causal
          influence} on $Y$;
  \end{itemize}
        plus a \textbf{conditional probability table (CPT)} at every node giving
        $P(\textit{node} \mid \textit{its parents})$.
  \item A root node has no parents, so its CPT is just its prior.
  \item \mk{The whole network in one equation}, and this is what the
        numericals use:
\end{itemize}
\[ P(x_1, \ldots, x_n) \;=\; \prod_{i=1}^{n} P\big(x_i \mid parents(x_i)\big) \]
\begin{itemize}
  \item \textbf{Why it is worth having}: a full joint distribution over $n$ binary
        variables needs $2^n - 1$ numbers. The alarm network below has 5 variables and
        needs \mk{only 10}, because each node depends on its parents
        alone. That compactness is the answer to \emph{when do we use a Bayesian
        network}.
  \item \textbf{Conditional independence} is the property being exploited: a node is
        independent of its non-descendants given its parents.
\end{itemize}}{%
\lead{\mk{Reasoning in the network} (the 5 or 6 marks)}
\begin{itemize}
  \item \textbf{Causal / top-down} $-$ from cause to effect. \emph{Burglary happened;
        how likely is John to call?} Uses the CPTs directly.
  \item \textbf{Diagnostic / bottom-up} $-$ from effect to cause. \emph{John called; how
        likely is a burglary?} This is where \textbf{Bayes' theorem enters}, and it is
        the answer to \emph{``how is Bayes' theorem used in a belief network''}
        \yr{(\textbf{77 Ch})}.
  \item \textbf{Intercausal} $-$ between two causes of one observed effect.
        \emph{The alarm rang and an earthquake is reported, so burglary becomes less
        likely}: \mk{explaining away}. Worth naming; almost nobody does.
  \item \textbf{Mixed} $-$ evidence above and below the query node at once.
  \item \textbf{The method for a numerical}: write the product form, substitute one CPT
        entry per variable, and read $P(\neg X) = 1 - P(X)$ off the same table.
\end{itemize}}

\penalty0\vspace{2pt}

\sbsr{0.30}{%
\figT{c4_causal_net.png}
{\color{sub}\footnotesize A causal net: \emph{Cloudy} influences \emph{Sprinkler} and
\emph{Rain}, both influence \emph{WetGrass}. The shape 78 Ba prints.}}{%
\figT{c4_bbn_alarm.png}
{\color{sub}\footnotesize The standard alarm network, with a CPT at every node.
\emph{Burglary} and \emph{Earthquake} both cause \emph{Alarm}; the alarm makes
\emph{John} and \emph{Mary} call. Five variables, and the graph plus these five tables
hold the whole joint distribution in \textbf{10 numbers} instead of $2^5 - 1 = 31$.
The worked line beside it is the product form applied straight down the network, which
is exactly what \yr{78 Ba} and \yr{\textbf{\texttt{82 Bh}}} ask for.}}

\penalty0\vspace{2pt}

\begin{itemize}
  \item \mk{The connection the papers ask for in one sentence}: a
        belief network stores \emph{causal} conditional probabilities, and Bayes' theorem
        is what lets you run them \emph{backwards} to a diagnosis. Without the network
        you would need the full joint distribution; without Bayes you could only reason
        in the direction the arrows point.
\end{itemize}
